\documentclass[11pt,a4paper]{article}
\usepackage{amsmath,amssymb}
\usepackage{fontspec}
\usepackage{xeCJK}
\setCJKmainfont{Hiragino Mincho ProN}
\setCJKsansfont{Hiragino Kaku Gothic ProN}
\setCJKmonofont{Osaka-Mono}
\usepackage{algorithm}
\usepackage{algpseudocode}
\usepackage{listings}
\usepackage{xcolor}
\usepackage{graphicx}
\usepackage{booktabs}
\usepackage{hyperref}
\usepackage[margin=2.5cm]{geometry}

% Code listing style
\lstset{
  language=C,
  basicstyle=\ttfamily\small,
  keywordstyle=\color{blue}\bfseries,
  commentstyle=\color{gray},
  stringstyle=\color{red},
  numbers=left,
  numberstyle=\tiny\color{gray},
  frame=single,
  breaklines=true,
  captionpos=b
}

\title{HPhi MPI通信バッチ処理最適化\\
\large 技術レポート}
\author{HPhi Development Team}
\date{\today}

\begin{document}

\maketitle

\begin{abstract}
本レポートでは、HPhi量子格子モデルシミュレータにおけるMPI通信の最適化手法について説明する。
複数の転送項やInterAll相互作用項が同じMPI通信相手（origin）を持つ場合、
これらをグループ化してMPI通信回数を削減するバッチ処理を実装した。
本手法により、特に大規模系・多プロセス環境でのMPI通信オーバーヘッドを大幅に削減できる。
\end{abstract}

\tableofcontents
\newpage

%==============================================================================
\section{はじめに}
%==============================================================================

HPhiは量子格子モデルの厳密対角化を行うシミュレータである。
大規模系の計算ではMPIによる並列化が必須であり、
ハミルトニアンの行列ベクトル積 $\hat{H}|\psi\rangle$ の計算において
プロセス間通信が発生する。

本最適化では、以下の問題に対処する：
\begin{itemize}
  \item 各転送項が個別にMPI\_Sendrecvを呼び出すことによる通信オーバーヘッド
  \item 同じ通信相手への複数回の重複した通信
\end{itemize}

%==============================================================================
\section{問題の背景}
%==============================================================================

\subsection{HPhiにおけるMPI並列化}

HPhiでは、ヒルベルト空間の基底状態を各MPIプロセスに分散して格納する。
$N$サイト系で$N_{\rm proc}$プロセスを使用する場合、
各プロセスは $N_{\rm site} = N - \log_2(N_{\rm proc})$ 個のサイトを
「ローカルサイト」として管理し、残りを「プロセス間サイト」として扱う。

\subsection{転送項のMPI通信}

ホッピング項 $c^\dagger_i c_j$ において、サイト$i$がローカル、
サイト$j$がプロセス間の場合（MPIsingle）、通信相手（origin）は
以下で計算される：
\begin{equation}
  \text{origin} = \text{myrank} \oplus \text{mask}
\end{equation}
ここで $\oplus$ はXOR演算、mask は $T_{\rm pow}[j]$（サイト$j$のビットマスク）である。

\subsection{従来実装の問題点}

従来の実装では、各転送項が個別にMPI\_Sendrecvを呼び出していた：

\begin{lstlisting}[caption=従来のMPIsingle処理（概念図）]
for (i = 0; i < N_transfers; i++) {
    origin = compute_origin(transfer[i]);
    MPI_Sendrecv(v1, ..., origin, ...);  // 毎回通信
    process_transfer(transfer[i], v1buf);
}
\end{lstlisting}

例えば、同じoriginへの転送が$N$個ある場合、$N$回のMPI\_Sendrecvが発生する。
これはMPI通信のレイテンシが支配的な場合に大きなオーバーヘッドとなる。

%==============================================================================
\section{最適化手法：バッチ処理}
%==============================================================================

\subsection{基本アイデア}

同じoriginを持つ転送項をグループ化し、1回のMPI\_Sendrecvで
必要なデータを取得した後、全ての転送項をローカルで処理する。

\begin{figure}[h]
\centering
\begin{tabular}{cc}
\textbf{最適化前} & \textbf{最適化後} \\
\hline
転送1 $\rightarrow$ MPI通信 & \\
転送2 $\rightarrow$ MPI通信 & グループ化 $\rightarrow$ 1回のMPI通信 \\
転送3 $\rightarrow$ MPI通信 & $\rightarrow$ 転送1,2,3をローカル処理 \\
\end{tabular}
\caption{バッチ処理の概念図（同じoriginへの3転送の場合）}
\end{figure}

\subsection{アルゴリズム}

\begin{algorithm}
\caption{MPIバッチ処理アルゴリズム}
\begin{algorithmic}[1]
\State \textbf{初期化フェーズ（1回のみ実行）:}
\For{各転送項 $t$}
    \State origin $\leftarrow$ myrank $\oplus$ mask[$t$]
    \State 転送$t$をorigin別にグループ化
    \State 係数・ローカルマスク等を事前計算
\EndFor
\State
\State \textbf{実行フェーズ（毎イテレーション）:}
\For{各グループ $g$}
    \State MPI\_Sendrecv(v1, origin[$g$], v1buf) \Comment{1回の通信}
    \For{グループ$g$内の各転送 $t$}
        \State v1bufを使って転送$t$を計算 \Comment{ローカル処理}
    \EndFor
\EndFor
\end{algorithmic}
\end{algorithm}

%==============================================================================
\section{データ構造}
%==============================================================================

\subsection{転送項グループ（MPITransferGroup）}

\begin{lstlisting}[caption=MPITransferGroup構造体]
typedef struct {
    int origin;                      // 通信相手のMPIランク
    int mask;                        // ビットマスク
    int num_transfers;               // このグループの転送数
    int *transfer_indices;           // 転送インデックス配列
    double complex *coefficients;    // 事前計算した係数
    unsigned long int *local_mask;   // ローカルサイトマスク
    unsigned long int *state1check;  // 状態チェック用マスク
    unsigned long int *bit1diff;     // フェルミオン符号計算用
} MPITransferGroup;
\end{lstlisting}

\subsection{バッチコンテナ（MPIBatchedTransfers）}

\begin{lstlisting}[caption=MPIBatchedTransfers構造体]
typedef struct {
    int num_groups;                  // グループ数（ユニークorigin数）
    MPITransferGroup *groups;        // グループ配列
    int is_initialized;              // 初期化フラグ
} MPIBatchedTransfers;
\end{lstlisting}

\subsection{MPIdouble用グループ（MPIDoubleTransferGroup）}

MPIdoubleでは両サイトがプロセス間のため、ローカルサイト情報は不要：

\begin{lstlisting}[caption=MPIDoubleTransferGroup構造体]
typedef struct {
    int origin;                      // 通信相手のMPIランク
    int num_transfers;               // このグループの転送数
    int *transfer_indices;           // 転送インデックス配列
    double complex *coefficients;    // 事前計算した係数（符号含む）
} MPIDoubleTransferGroup;

typedef struct {
    int num_groups;                  // グループ数
    MPIDoubleTransferGroup *groups;  // グループ配列
    int is_initialized;              // 初期化フラグ
} MPIBatchedDoubleTransfers;
\end{lstlisting}

\subsection{InterAll用グループ（MPIInterAllGroup）}

4体相互作用 $c^\dagger_i c_j c^\dagger_k c_l$ 用の構造体：

\begin{lstlisting}[caption=MPIInterAllGroup構造体]
typedef struct {
    int origin;                      // 通信相手
    int num_interall;                // InterAll項数
    int *interall_indices;           // インデックス配列
    double complex *coefficients;    // 係数
    int *is_hermite;                 // エルミート共役フラグ
    unsigned long int *tmp_isite1;   // 事前計算サイト情報
    unsigned long int *tmp_isite2;
    unsigned long int *tmp_isite3;
    unsigned long int *tmp_isite4;
    // ... (Tpow値も格納)
} MPIInterAllGroup;
\end{lstlisting}

%==============================================================================
\section{実装詳細}
%==============================================================================

\subsection{対応モデル}

\begin{table}[h]
\centering
\begin{tabular}{lccc}
\toprule
モデル & MPIsingleバッチ & MPIdoubleバッチ & InterAllバッチ \\
\midrule
SpinlessFermionGC & $\checkmark$ & -- & -- \\
SpinlessFermion (canonical) & $\checkmark$ & -- & -- \\
HubbardGC & $\checkmark$ & $\checkmark$ & $\checkmark$ \\
Hubbard (canonical) & $\checkmark$ & $\checkmark$ & -- \\
SpinGC & $\checkmark$ (Exchange) & -- & -- \\
Spin (canonical) & $\checkmark$ (Exchange) & -- & -- \\
\bottomrule
\end{tabular}
\caption{バッチ処理対応状況}
\end{table}

\subsection{MPIsingle vs MPIdouble}

\begin{itemize}
  \item \textbf{MPIsingle}: 一方のサイトがローカル、他方がプロセス間
  \item \textbf{MPIdouble}: 両方のサイトがプロセス間
\end{itemize}

MPIdoubleはMPIsingleよりも単純で、受信後の処理が単なるベクトルスケーリングになる：
\begin{equation}
  v_0[j] \mathrel{+}= \text{trans} \times v_{\rm buf}[j]
\end{equation}

\subsection{MPIdoubleバッチ処理の詳細}

\subsubsection{origin計算（MPIdouble）}

両サイトがプロセス間の場合、originは以下で計算される：
\begin{equation}
  \text{origin} = \text{myrank} \oplus (\text{mask}_1 + \text{mask}_2)
\end{equation}
ここで $\text{mask}_1 = T_{\rm pow}[2 \cdot \text{site}_1 + \sigma_1]$、
$\text{mask}_2 = T_{\rm pow}[2 \cdot \text{site}_2 + \sigma_2]$ である。

\subsubsection{フェルミオン符号（MPIdouble）}

MPIdoubleでのフェルミオン符号は以下で計算される：
\begin{lstlisting}[caption=MPIdoubleフェルミオン符号の計算]
if (mask2 > mask1) bitdiff = mask2 - mask1 * 2;
else bitdiff = mask1 - mask2 * 2;
SgnBit((unsigned long int)(origin & bitdiff), &Fsgn);
\end{lstlisting}

\subsubsection{状態の有効性チェック}

各プロセスで転送が有効かどうかは、originの状態で判定される：
\begin{itemize}
  \item $\text{state}_1 = 0$ かつ $\text{state}_2 = \text{mask}_2$:
        順方向（$c^\dagger_1 c_2$）、係数 $= -F_{\rm sgn} \times t$
  \item $\text{state}_1 = \text{mask}_1$ かつ $\text{state}_2 = 0$:
        逆方向（$c^\dagger_2 c_1$）、係数 $= -F_{\rm sgn} \times t^*$
  \item その他: 無効状態、係数 $= 0$
\end{itemize}

\textbf{重要}: 全プロセスが同じ通信パターンを持つ必要があるため、
無効状態でもMPI通信は実行し、係数を0にすることで寄与をなくす。

\subsubsection{HubbardGC vs Hubbard（canonical）}

\begin{table}[h]
\centering
\begin{tabular}{lcc}
\toprule
項目 & HubbardGC & Hubbard (canonical) \\
\midrule
list\_1交換 & 不要 & 必要 \\
インデックス計算 & 直接 $j$ & GetOffComp(list\_1buf[$j$]) \\
処理の複雑さ & 単純 & やや複雑 \\
\bottomrule
\end{tabular}
\caption{HubbardGCとHubbard（canonical）のMPIdouble処理の違い}
\end{table}

Hubbard（canonical）では、粒子数保存のため状態空間が制限される。
そのため、list\_1（状態インデックス配列）を交換し、
GetOffComp関数でターゲットインデックスを計算する必要がある：

\begin{lstlisting}[caption=Hubbard canonical MPIdoubleの処理]
// list_1の交換
MPI_Sendrecv(list_1, ..., list_1buf, ...);
MPI_Sendrecv(tmp_v1, ..., v1buf, ...);

// 各要素の処理
for (j = 1; j <= idim_max_buf; j++) {
    GetOffComp(list_2_1, list_2_2, list_1buf[j],
               irght, ilft, ihfbit, &ioff);
    tmp_v0[ioff] += trans * v1buf[j];
}
\end{lstlisting}

\subsection{Spinモデルのバッチ処理}

\subsubsection{SpinモデルのExchange相互作用}

Spin 1/2モデルでは、Exchange相互作用 $J \mathbf{S}_i \cdot \mathbf{S}_j$
のスピン反転項がMPI通信を必要とする。
Hubbardモデルとの主な違いは、ビットマスクの計算方法である：

\begin{itemize}
  \item \textbf{Hubbard}: mask $= T_{\rm pow}[2 \cdot \text{site} + \sigma]$
  \item \textbf{Spin}: mask $= T_{\rm pow}[\text{site}]$
\end{itemize}

Spinモデルではスピン自由度が1サイトあたり1ビットであるため、
サイト番号がそのままビット位置に対応する。

\subsubsection{Spinモデル用グループ構造体}

\begin{lstlisting}[caption=MPISpinExchangeGroup構造体]
typedef struct {
    int origin;                      // 通信相手のMPIランク
    int num_terms;                   // Exchange項数
    int *term_indices;               // インデックス配列
    double complex *coefficients;    // 係数
    int *org_isite1;                 // ローカルサイト
    int *org_ispin1;                 // フリップ前スピン
    int *org_ispin2;                 // フリップ後スピン
    int *state1check;                // 状態チェック値
} MPISpinExchangeGroup;
\end{lstlisting}

\subsubsection{SpinGC vs Spin（canonical）}

\begin{table}[h]
\centering
\begin{tabular}{lcc}
\toprule
項目 & SpinGC & Spin (canonical) \\
\midrule
list\_1交換 & 不要 & 必要 \\
処理関数 & X\_child\_GC\_CisAitCiuAiv\_spin\_MPIsingle\_batched &
           X\_child\_general\_int\_spin\_MPIsingle\_batched \\
インデックス計算 & 直接 & GetOffComp使用 \\
\bottomrule
\end{tabular}
\caption{SpinGCとSpin（canonical）のMPIsingle処理の違い}
\end{table}

\subsubsection{Spinモデルテスト結果}

8サイト鎖のHeisenbergモデル（4 MPIプロセス）でのテスト：

\begin{table}[h]
\centering
\begin{tabular}{lccc}
\toprule
モデル & Exchange項数 & グループ数 & エネルギー \\
\midrule
SpinGC & 8 & 2 & $-3.6510934089$ \\
Spin (canonical, $S_z=0$) & 8 & 2 & $-3.6510934089$ \\
\bottomrule
\end{tabular}
\caption{Spinモデルバッチ処理のテスト結果}
\end{table}

両モデルで同じ基底状態エネルギーが得られ、バッチ処理の正確性を確認した。

\subsection{origin計算}

\subsubsection{転送項（Hubbard/SpinlessFermion）}

サイト$i$（ローカル）からサイト$j$（プロセス間）への転送：
\begin{equation}
  \text{origin} = \text{myrank} \oplus T_{\rm pow}[2j + \sigma]
\end{equation}
ここで $\sigma$ はスピン（Hubbardモデル）、SpinlessFermionでは $\sigma=0$。

\subsubsection{InterAll項}

4サイト相互作用では、CheckBit\_InterAllPE関数により
ビット操作を適用してoriginを計算：
\begin{equation}
  \text{origin} = A^\dagger_1 C_2 A^\dagger_3 C_4 |\text{myrank}\rangle
\end{equation}

\subsection{フェルミオン符号}

フェルミオン系では、演算子の交換により符号が変化する。
バッチ処理では初期化時に符号を事前計算：

\begin{lstlisting}[caption=フェルミオン符号の事前計算]
// プロセス間部分の符号
int bit2diff = mask2 - 1;
int Fsgn;
SgnBit((unsigned long int)(origin & bit2diff), &Fsgn);

// 係数に符号を含める
coefficients[t] = -(double)Fsgn * trans_coeff;
\end{lstlisting}

%==============================================================================
\section{性能評価}
%==============================================================================

\subsection{削減効果（MPIsingle）}

4$\times$2格子のHubbardGCモデル（8サイト、4 MPIプロセス）での測定結果：

\begin{table}[h]
\centering
\begin{tabular}{lcc}
\toprule
項目 & 最適化前 & 最適化後 \\
\midrule
MPIsingle転送数 & 6 & 6 \\
MPI\_Sendrecv回数 & 6 & 2 \\
削減率 & -- & \textbf{3.0倍} \\
\bottomrule
\end{tabular}
\caption{MPI通信回数の削減（MPIsingle）}
\end{table}

\subsection{MPIdoubleテスト結果}

4サイト鎖のHubbardモデル（16 MPIプロセス）でのMPIdoubleバッチ処理テスト：

\begin{table}[h]
\centering
\begin{tabular}{lccc}
\toprule
モデル & MPIdouble転送数 & グループ数 & エネルギー \\
\midrule
HubbardGC & 2 & 2 & $-3.41855071887385$ \\
Hubbard (canonical) & 2 & 2 & $-2.10274848346207$ \\
\bottomrule
\end{tabular}
\caption{MPIdoubleバッチ処理のテスト結果}
\end{table}

16 MPIプロセスを使用することで、サイト2,3がプロセス間サイトとなり、
これらのサイト間の転送がMPIdoubleとして処理される。

\subsection{デバッグ出力}

実行時に以下のようなバッチ処理統計が出力される：
\begin{verbatim}
[MPI Batching] HubbardGC MPIdouble: 2 transfers -> 2 groups (1.0x reduction)
[MPI Batching] HubbardGC: 4 MPIsingle transfers -> 4 groups (1.0x reduction)
[MPI Batching] HubbardGC InterAll: No off-diagonal terms
\end{verbatim}

Hubbard（canonical）の場合：
\begin{verbatim}
[MPI Batching] Hubbard MPIdouble: 2 transfers -> 2 groups (1.0x reduction)
[MPI Batching] HubbardGC: 4 MPIsingle transfers -> 4 groups (1.0x reduction)
\end{verbatim}

\subsection{期待される効果}

\begin{itemize}
  \item \textbf{最良ケース}: 同じoriginへの$N$転送 $\rightarrow$ MPI呼び出し$N$回が1回に
  \item \textbf{典型的改善}: 2D/3Dモデルで顕著（多くの転送が同じoriginを共有）
  \item \textbf{オーバーヘッド}: 初期化時に1回のグループ化コスト（無視できる程度）
\end{itemize}

%==============================================================================
\section{使用方法}
%==============================================================================

バッチ処理は自動的に有効化される。特別な設定は不要である。

\subsection{ビルド}

\begin{verbatim}
mkdir build && cd build
cmake .. -DUSE_MPI=ON
make
\end{verbatim}

\subsection{実行}

通常通りMPIで実行：
\begin{verbatim}
mpirun -np 4 ./HPhi -s stan.in
\end{verbatim}

%==============================================================================
\section{ソースコード構成}
%==============================================================================

\begin{table}[h]
\centering
\begin{tabular}{ll}
\toprule
ファイル & 内容 \\
\midrule
\texttt{src/include/mltplyMPIBatched.h} & データ構造・関数宣言 \\
\texttt{src/mltplyMPIBatched.c} & バッチ処理実装 \\
\texttt{src/mltplyHubbard.c} & Hubbardモデル（バッチ呼び出し） \\
\texttt{src/mltplySpinless.c} & SpinlessFermionモデル（バッチ呼び出し） \\
\texttt{src/mltplySpin.c} & Spinモデル（バッチ呼び出し） \\
\bottomrule
\end{tabular}
\caption{関連ソースファイル}
\end{table}

%==============================================================================
\section{まとめ}
%==============================================================================

本最適化により、HPhiのMPI通信効率を改善した。
同じ通信相手を持つ複数の転送項やInterAll項をグループ化し、
MPI\_Sendrecv呼び出し回数を削減することで、
特に通信レイテンシが支配的な環境での性能向上が期待できる。

\subsection{実装済み機能}

\begin{itemize}
  \item MPIsingleバッチ処理:
    \begin{itemize}
      \item SpinlessFermionGC, SpinlessFermion
      \item HubbardGC, Hubbard
      \item SpinGC, Spin（Exchange相互作用）
    \end{itemize}
  \item MPIdoubleバッチ処理: HubbardGC, Hubbard（canonical）
  \item InterAllバッチ処理: HubbardGC
\end{itemize}

\subsection{今後の拡張}

今後の拡張として以下が考えられる：
\begin{itemize}
  \item Hubbard (canonical) のInterAllバッチ処理
  \item SpinモデルのMPIdoubleバッチ処理
  \item SpinlessFermionのMPIdoubleバッチ処理
  \item GeneralSpinモデルへの適用
\end{itemize}

%==============================================================================
\section*{謝辞}
%==============================================================================

本実装はClaude Opus 4.5との協働により開発された。

\end{document}
